\chapter{Kết luận và hướng phát triển}
\label{chap:conclusion}

Đồ án đã hoàn thành mục tiêu nghiên cứu về bài toán chọn tập chuỗi protein đại diện có ưu tiên nhóm hiếm thông qua mô hình tối ưu hóa submodular có trọng số. Từ việc chỉ ra hạn chế phương pháp luận của hàm mục tiêu không trọng số truyền thống tại chương~\ref{chap:theory}, thiết lập khung phương pháp luận RareProtSelect và chứng minh chặt chẽ các bảo đảm toán học tại chương~\ref{chap:methodology}, đến việc thực thi chuỗi 6 thử nghiệm tuần tự có kiểm soát tại chương~\ref{chap:experiments}, đồ án đã khẳng định tính đúng đắn cả về mặt lý thuyết giải tích lẫn thực nghiệm tính toán sinh học. Chương~\ref{chap:conclusion} hệ thống hóa các kết quả và đóng góp khoa học chính, làm rõ ranh giới kiểm chứng cùng các hạn chế nội tại, đồng thời mở ra các định hướng nghiên cứu tiếp theo nhằm hoàn thiện và mở rộng khung phương pháp.

\section{Tổng kết đóng góp khoa học và kết luận định hướng ứng dụng}
\label{sec:conc_contributions}

Thay vì liệt kê lặp lại số liệu của từng thử nghiệm riêng lẻ đã được tổng hợp tại bảng~\ref{tab:rq_resolution}, phần này đúc kết các phát hiện khoa học ở cấp độ phương pháp luận: đánh giá bản chất vì sao khung phương pháp đề xuất thực sự hiệu quả, làm sáng tỏ vai trò củng cố của các khảo sát mở rộng đa chiều, và xác lập hướng dẫn lựa chọn cấu hình cụ thể cho từng tình huống sinh học thực tế.

\subsection{Đánh giá tính hiệu quả của phương pháp đề xuất và cơ chế giải quyết đánh đổi đa mục tiêu}
\label{subsec:conc_core_effectiveness}
\phantomsection\label{subsec:conc_five_tasks}
\phantomsection\label{subsec:conc_core_methodology}

Đồ án đã giải quyết đầy đủ năm nhiệm vụ nghiên cứu đặt ra tại phần Mở đầu: từ việc sàng lọc và thiết lập tập dữ liệu phân họ chính thức $V_{\text{curated}}$ gồm $4.120$ chuỗi thuộc $115$ Subfamily từ $13.847$ chuỗi thô của tệp MEROPS\_A1A.fasta~\cite{merops} tại mục~\ref{subsec:methodology_curation}; tái lập độc lập mô hình cơ sở của Nguyen và cộng sự~\cite{nguyen2025} trên $13.307$ chuỗi với sai số tuyệt đối trung bình chỉ $0{,}51\%$ tại bảng~\ref{tab:exp1_milestones} và vạch rõ sai lệch bao phủ chéo ngoài ranh giới phân họ; thiết kế họ hàm mục tiêu Weighted Facility Location $F_\alpha(S)$ tại công thức~\eqref{eq:sample_weight_formula}--\eqref{eq:weighted_facility_location}, chứng minh tính đơn điệu không giảm tại bổ đề~\ref{lem:monotonicity}, tính submodular tại định lý~\ref{thm:submodularity}, thiết lập nền tảng lý thuyết theo cận xấp xỉ Nemhauser $(1 - 1/e) \approx 63{,}2\%$ của giải thuật tham lam chuẩn tại hệ quả~\ref{cor:nemhauser_guarantee}~\cite{nemhauser1978} làm cơ sở cho giải thuật thực thi Pruned Greedy 90\% trên GPU tại thuật toán~\ref{alg:gpu_pruned_greedy}; đến việc khảo sát không gian ngữ nghĩa ESM-2~\cite{esm2} kết hợp dịch tâm Mean-Centering với $\beta = 0{,}70$ tại thuật toán~\ref{alg:mean_centering}; và hoàn tất kiểm chứng chéo hai chiều, phân tích phân bổ vi mô cùng khảo sát mở rộng ngân sách. Từ toàn bộ chuỗi nghiên cứu này, kết luận cốt lõi được khẳng định là phương pháp đề xuất RareProtSelect mang lại hiệu quả thực chất nhờ giải quyết đồng thời hai mục tiêu: vừa khôi phục khả năng nhận diện các Subfamily hiếm, vừa duy trì độ bao phủ toàn cục, thay vì phải đánh đổi hay hy sinh hiệu năng tổng thể như các cách tiếp cận thông thường.

Trong các nghiên cứu chọn tập đại diện truyền thống, việc gia tăng tính đại diện cho các nhóm hiếm thường bị mặc định là mục tiêu đối nghịch trực tiếp với việc tối đa hóa độ bao phủ toàn cục. Đồ án đã chỉ ra rằng tại mốc ngân sách đủ lớn trên các siêu họ protein có phân bố đuôi dài với tỷ lệ chênh lệch quy mô gấp $329{,}0$ lần theo công thức~\eqref{eq:imbalance_ratio}, hai mục tiêu này có thể đồng thời đạt được. Khi áp dụng thước đo nội phân họ Within-Subfamily Coverage theo công thức~\eqref{eq:within_subfamily_reps}--\eqref{eq:within_subfamily_covered_set}, mô hình không trọng số ($\alpha = 0{,}0$) tại ngân sách chuẩn $k = 115$ bộc lộ sự lãng phí ngân sách đại diện: giải thuật dồn tới $58$ vị trí đại diện trùng lặp vào các Subfamily lớn đã bão hòa độ bao phủ nội bộ, trong khi bỏ sót $58$ trên $115$ Subfamily, chiếm $50{,}43\%$, ở mức $0{,}0\%$ Sensitivity. Bằng cách điều chỉnh tổng trọng số của mỗi phân họ theo hàm lũy thừa $W_g^{(\alpha)} = (n_g)^{1-\alpha}$ tại công thức~\eqref{eq:subfamily_total_weight}, cấu hình trọng số cân bằng $\alpha = 0{,}5$ thu hẹp tỷ lệ chênh lệch đóng góp giữa Subfamily lớn nhất A01.009 và Subfamily đơn độc từ $329{,}0$ lần xuống $18{,}14$ lần. Cơ chế này không làm mất đi các chuỗi đại diện thiết yếu của các Subfamily lớn, mà chỉ hạn chế các lượt chọn trùng lặp có độ tăng cận biên thấp ở nhóm lớn để chuyển dịch ngân sách sang các phân họ chưa có đại diện. Cần nhấn mạnh rằng sự giải tỏa đánh đổi đa mục tiêu này đạt được tại mốc ngân sách đủ lớn $k = G = 115$ với cấu hình cân bằng $\alpha = 0{,}5$, nơi các Subfamily lớn đã đạt bão hòa nội nhóm; trong khi ở các mức ngân sách nhỏ và vừa ($k < G$), việc bảo tồn nhóm hiếm vẫn đòi hỏi sự đánh đổi một phần chỉ số Overall Sensitivity, chẳng hạn tại $k = 30$, cấu hình $\alpha = 0{,}5$ đạt $4{,}44\%$ Rare Sensitivity nhưng chỉ số Overall Sensitivity giảm $8{,}93$ điểm phần trăm từ $71{,}65\%$ xuống $62{,}72\%$ so với $\alpha = 0{,}0$. Tại $k = 115$, trên không gian NCBI BLASTp~\cite{blast}, cấu hình $\alpha = 0{,}5$ nâng chỉ số Rare Sensitivity từ $2{,}22\%$ lên $59{,}11\%$, tức tăng gấp $26{,}6$ lần khi cố định ngưỡng ngân sách tối đa $k_{\max} = 115$, và đạt $50{,}22\%$, tức tăng gấp $22{,}6$ lần khi ngưỡng ngân sách tối đa mở rộng lên $k_{\max} = 380$; đồng thời chỉ số Macro-Average Sensitivity tăng từ $49{,}28\%$ lên $78{,}87\%$ trong khi chỉ số Overall Sensitivity đạt $94{,}13\%$ so với $93{,}57\%$ của mô hình cơ sở. Đây là bằng chứng thực nghiệm khẳng định tính hiệu quả có điều kiện biên của phương pháp đề xuất.

\subsection{Giá trị củng cố luận điểm khoa học từ các khảo sát mở rộng đa chiều}
\label{subsec:conc_reinforcing_surveys}
\phantomsection\label{subsec:conc_exploratory_esm2}

Bên cạnh kết quả trọng tâm trên không gian tương đồng chuỗi BLASTp ở ngân sách cơ sở, bốn thử nghiệm mở rộng từ thử nghiệm 3 đến thử nghiệm 6 tại chương~\ref{chap:experiments} hợp thành ba tầng bằng chứng độc lập làm rõ tính nhất quán, tính tường minh của cơ chế tối ưu hóa và khả năng mở rộng quy mô của phương pháp đề xuất:

\begin{enumerate}
    \item \textbf{Củng cố tính tổng quát không phụ thuộc không gian biểu diễn và quy luật tương tác cấu trúc ma trận:} Việc tích hợp module khảo sát mở rộng ESM-2 kết hợp kỹ thuật dịch tâm Mean-Centering với $\beta = 0{,}70$ tại thử nghiệm 3 và bảng~\ref{tab:exp3_milestones} chứng minh công thức trọng số nghịch đảo quy mô không bị phụ thuộc cục bộ vào đặc thù ma trận thưa BLASTp. Ngay trên không gian hình học liên tục có cấu trúc đặc với $99{,}64\%$ phần tử dương sau khi dịch tâm, trọng số $\alpha = 0{,}5$ vẫn nhất quán nâng chỉ số Rare Coverage ESM tăng gấp $3{,}4$ lần, từ $12{,}85\%$ lên $43{,}21\%$. Đặc biệt, ma trận kiểm chứng chéo hai chiều $2 \times 2$ trên $120$ cấu hình tại thử nghiệm 4 và bảng~\ref{tab:exp4_cross_milestones} đã củng cố hiểu biết lý thuyết về quy luật tương tác giữa cấu trúc ma trận và quy mô ngân sách $k$: ma trận đặc ESM-2 phát huy ưu thế dẫn truyền thông tin ngữ nghĩa liên cụm khi ngân sách hẹp ($k \in [8..75]$), giúp nhận diện sớm nhóm hiếm ngay từ $k = 8$ và đạt hiệu năng cao hơn gấp $1{,}6$ đến $4{,}0$ lần so với BLASTp; ngược lại, tính cô lập phân họ của ma trận thưa BLASTp giúp duy trì độ tăng cận biên cho các phân họ hiếm còn sót lại khi ngân sách đủ lớn ($k = 115$), đạt tỷ lệ bao phủ $92$ trên $115$ Subfamily so với $77$ trên $115$ của ESM-2.
    
    \item \textbf{Làm sáng tỏ cơ chế phân bổ đại diện bên trong quá trình tối ưu hóa qua phân tích vi mô:} Khảo sát cấu trúc nhãn phân họ trực tiếp của từng chuỗi trong tập nghiệm $S_k$ tại thử nghiệm 5 và bảng~\ref{tab:exp5_milestones} cung cấp bằng chứng nhân quả giải thích trực tiếp kết quả của thử nghiệm 2 và thử nghiệm 3. Số liệu chứng minh sự cải thiện tỷ lệ nhận diện chuỗi không bắt nguồn từ nhiễu ngẫu nhiên, mà do trọng số $\alpha = 0{,}5$ đã cắt giảm độ dư thừa phân bổ Subfamily Redundancy từ $50{,}43\%$ xuống $20{,}00\%$, qua đó chuyển dịch đúng $35$ vị trí đại diện từng bị phân bổ trùng lặp ở các nhóm lớn, điển hình là giảm một nửa tỷ trọng của Subfamily lớn nhất A01.009 tại mốc chuẩn $k = 115$ từ $6$ vị trí ($5{,}22\%$) ở $\alpha = 0{,}0$ xuống còn $3$ vị trí ($2{,}61\%$) ở $\alpha = 0{,}5$ (và tương tự tại $k = 11$ giảm từ $18{,}18\%$ xuống $9{,}09\%$), sang cho $35$ Subfamily chưa có đại diện. Điều này giúp nâng số lượng phân họ hiếm có đại diện từ $1$ lên $27$ trên $45$ Subfamily, đạt $60{,}00\%$ ở $\alpha = 0{,}5$ và đạt $41$ trên $45$ Subfamily, tức $91{,}11\%$ ở $\alpha = 1{,}0$.
    
    \item \textbf{Củng cố năng lực bao phủ đồng thời nhóm hiếm và biến dị nội nhóm của các phân họ lớn khi mở rộng ngân sách:} Thực nghiệm mở rộng ngân sách theo cấu trúc phân tầng lên $k = 380$ tại thử nghiệm 6 và bảng~\ref{tab:exp6_milestones} giải đáp câu hỏi về việc giảm số đại diện của Subfamily lớn nhất A01.009 xuống $3$ chuỗi tại $k = 115$ liệu có làm mất thông tin biến dị nội nhóm của phân họ có quy mô lớn nhất này hay không. Phân họ A01.009 vốn đã đạt $100{,}0\%$ Sensitivity nội phân họ ngay tại $k = 115$; khi ngân sách được nới rộng lên $k = 380$, hàm mục tiêu $\alpha = 0{,}5$ không chỉ nâng độ bao phủ lên $112$ trên $115$ Subfamily, đạt $97{,}39\%$, với $104$ Subfamily đạt bão hòa $100\%$ Sensitivity, mà còn tự động phân bổ bổ sung thành $6$ chuỗi đại diện cho A01.009 trên không gian BLASTp (nâng Coverage ESM từ $91{,}88\%$ lên $93{,}67\%$) và $22$ chuỗi trên không gian ESM-2 (nâng Coverage ESM từ $93{,}27\%$ lên $95{,}81\%$), nâng cao độ bao phủ biến dị ngữ nghĩa liên tục cho phân họ lớn này.
\end{enumerate}

\subsection{Khuyến nghị lựa chọn cấu hình theo tình huống thực tế}
\label{subsec:conc_config_guidelines}

Từ các bằng chứng thực nghiệm và phân tích cơ chế trên cả hai không gian biểu diễn, đồ án xác lập bộ khuyến nghị lựa chọn cấu hình RareProtSelect cho bốn tình huống ứng dụng sinh học:

\begin{enumerate}
    \item \textbf{Tình huống 1: Chọn tập đại diện chuẩn đa mục tiêu ở ngân sách cơ sở tương đương số lượng phân họ ($k \approx G$).} Khuyến nghị sử dụng không gian tương đồng chuỗi NCBI BLASTp kết hợp hệ số trọng số cân bằng $\alpha = 0{,}5$ tại ngân sách $k = G$, tương ứng $k = 115$ trên Family A1. Đây là cấu hình mặc định tối ưu cho đại đa số các bài toán tiền xử lý so khớp đa chuỗi, dựng cây phát sinh chủng loài toàn siêu họ và xây dựng tập dữ liệu huấn luyện cân bằng. Ở cấu hình này, người dùng đạt được trạng thái cân bằng đa mục tiêu tối ưu: duy trì chỉ số Overall Sensitivity ở mức cao đạt $94{,}13\%$, khôi phục đáng kể tỷ lệ nhận diện phân họ hiếm đạt $59{,}11\%$, gấp $26{,}6$ lần mô hình cơ sở, đại diện cho $80{,}00\%$ tổng số Subfamily và giảm thiểu độ dư thừa xuống chỉ còn $20{,}00\%$.
    
    \item \textbf{Tình huống 2: Sàng lọc nhanh với ngân sách siêu tinh gọn ($k \ll G$) hoặc khảo sát các họ protein phân kỳ tiến hóa xa.} Khuyến nghị sử dụng không gian ngữ nghĩa liên tục ESM-2 kết hợp dịch tâm Mean-Centering với $\beta = 0{,}70$ và trọng số $\alpha = 0{,}5$ ở dải ngân sách hẹp $k \in [8..75]$. Cấu hình này phù hợp cho các bước phân tích sinh học tiếp theo có chi phí tính toán hoặc chi phí thực nghiệm cao trên mỗi chuỗi, chẳng hạn như mô phỏng động lực học phân tử, sàng lọc tương tác thuốc trên cấu trúc không gian ba chiều hoặc tổng hợp protein thực nghiệm trong phòng thí nghiệm, nơi số lượng chuỗi đại diện được phép chọn nhỏ hơn nhiều so với tổng số phân họ. Khi đó, khả năng nắm bắt tương đồng cấu trúc tầm xa của ma trận đặc ESM-2 giúp tập đại diện siêu nhỏ vẫn nhận diện sớm các nhóm hiếm ngay từ $k = 8$ và đạt hiệu năng bảo tồn nhóm hiếm cao hơn gấp $1{,}6$ đến $4{,}0$ lần so với sử dụng BLASTp cùng mức trọng số $\alpha = 0{,}5$. Tuy nhiên, nếu mục tiêu thực nghiệm tại $k = 60$ là tối đa hóa trực tiếp tỷ lệ nhận diện Live BLASTp Sensitivity, việc sử dụng BLASTp với trọng số mạnh $\alpha = 1{,}0$ đạt kết quả cao hơn rõ rệt, với $57{,}04\%$ Rare Sensitivity so với $28{,}89\%$ của ESM-2 $\alpha = 0{,}5$ và $40{,}44\%$ của ESM-2 $\alpha = 1{,}0$, trong khi vẫn duy trì chỉ số Overall Sensitivity ở mức $62{,}57\%$.
    
    \item \textbf{Tình huống 3: Nghiên cứu chuyên sâu về tiến hóa nhóm hiếm và bảo tồn tối đa các nhánh đơn độc.} Khuyến nghị sử dụng không gian tương đồng chuỗi NCBI BLASTp kết hợp hệ số trọng số bình đẳng nhóm $\alpha = 1{,}0$ tại ngân sách $k \ge G$. Cấu hình này phù hợp cho các nghiên cứu sinh học tiến hóa chuyên sâu nhằm lập bản đồ các nhánh phân loại hiếm, tìm kiếm các enzyme có cơ chế xúc tác đặc hữu hoặc bảo vệ các Subfamily đơn độc chỉ có duy nhất $1$ chuỗi trong cơ sở dữ liệu. Nhờ thiết lập tổng trọng số của mọi phân họ ngang nhau ($W_g^{(1{,}0)} = 1$), cấu hình này đẩy chỉ số Rare Sensitivity lên mức cao nhất $91{,}11\%$, bảo tồn $41$ trên $45$ Subfamily hiếm, tăng gấp $41{,}0$ lần so với mô hình không trọng số, phù hợp khi người nghiên cứu ưu tiên cao nhất độ bao phủ phân loại và chấp nhận giảm mật độ lấy mẫu ở các phân họ phổ biến, với Overall Sensitivity ở mức $79{,}08\%$.
    
    \item \textbf{Tình huống 4: Xây dựng thư viện tham chiếu đầy đủ cho các siêu họ có biến dị nội nhóm lớn.} Khuyến nghị sử dụng không gian tương đồng chuỗi NCBI BLASTp kết hợp hệ số trọng số cân bằng $\alpha = 0{,}5$ và mở rộng ngân sách theo cấu trúc phân tầng lên $k \approx 2{,}5G \text{ đến } 3{,}3G$, tương ứng $k \in [270..380]$ trên Family A1. Cấu hình này được khuyến nghị khi xây dựng bộ dữ liệu chuỗi đại diện tham chiếu hoặc bộ khung phân loại cho toàn siêu họ, nơi đòi hỏi cả hai mục tiêu đều phải tiệm cận mức bão hòa: vừa bao phủ $97{,}39\%$ tổng số phân họ, tức $112$ trên $115$ Subfamily, và đạt $93{,}33\%$ Rare Sensitivity, vừa phân bổ đủ từ $5$ đến $6$ chuỗi đại diện cho các phân họ quy mô lớn có biến dị nội nhóm rộng như A01.009 để duy trì $100{,}0\%$ Sensitivity và nâng Coverage ESM lên $93{,}67\%$, đồng thời giúp $104$ trên $115$ Subfamily đạt bão hòa $100\%$ Sensitivity.
\end{enumerate}

\section{Ranh giới kiểm chứng và các hạn chế tồn tại}
\label{sec:conc_limitations}

Bên cạnh những kết quả đạt được, việc xác định rõ phạm vi kiểm chứng thực nghiệm và các giới hạn nội tại của mô hình là cơ sở cần thiết để ứng dụng phương pháp một cách phù hợp và định hướng cải tiến trong tương lai. Qua quá trình thực nghiệm trên siêu họ Peptidase Family A1, khung phương pháp luận RareProtSelect còn tồn tại ba nhóm hạn chế phương pháp luận cần được phân tích làm rõ.

\subsection{Sự phụ thuộc vào nhãn phân họ có sẵn, sự gắn kết nhãn của thước đo và phạm vi các phương pháp đối chứng cố định ngân sách}
\label{subsec:lim_unassigned_upa}

Hạn chế căn bản đầu tiên nằm ở sự phụ thuộc vào hệ thống nhãn phân loại có sẵn. Công thức trọng số nghịch đảo quy mô $w_v^{(\alpha)} = (n_{g(v)})^{-\alpha}$ được xây dựng dựa trên tiền đề mỗi chuỗi protein $v \in V$ đều đã có nhãn phân họ xác định $g(v) \in \{1, \dots, G\}$. 

Đối với cơ sở dữ liệu MEROPS, tệp MEROPS\_A1A.fasta chứa tới $9.187$ chuỗi mang nhãn [A01.UPA], chiếm $69{,}04\%$ tổng số chuỗi đã qua kiểm soát chất lượng. Việc sàng lọc loại bỏ khối [A01.UPA] để xây dựng tập dữ liệu phân họ chính thức $V_{\text{curated}}$ gồm $4.120$ chuỗi là giải pháp bắt buộc nhằm bảo đảm tính chính xác cho phép đo nội phân họ. Tuy nhiên, điều này giới hạn phạm vi kiểm chứng hiện tại trong không gian các chuỗi đã có định danh phân loại, chưa thể trực tiếp xử lý các tập dữ liệu metagenomics môi trường chứa nhiều chuỗi chưa phân loại.

Đồng thời, việc cả trọng số tối ưu hóa $w_v^{(\alpha)}$ lẫn cơ chế đánh giá Within-Subfamily Coverage ($S_k \cap V_g \to V_g$) đều dựa trên cùng một hệ thống phân hoạch Subfamily tạo ra một sự tương thích có giám sát giữa mục tiêu huấn luyện và thước đo đánh giá. Khi so sánh với mô hình không trọng số ($\alpha = 0{,}0$) vốn không sử dụng thông tin nhãn, cấu hình $\alpha > 0$ hiển nhiên nắm giữ lợi thế về mặt thông tin cấu trúc. Hơn nữa, nếu nhãn phân họ đã biết trước tại đúng mốc ngân sách $k = G = 115$, một phương pháp đối chứng hạn mức cứng đơn giản là 1-medoid-per-subfamily có thể đạt $100\%$ Subfamily Representation và $100\%$ Rare Representation, cao hơn mức $80{,}00\%$ của $\alpha = 0{,}5$ và $84{,}35\%$ của $\alpha = 1{,}0$. Dù chiến lược 1-medoid-per-subfamily không giải quyết được bài toán xếp hạng liên tục ở $k \ll G$ hay phân bổ đa đại diện theo độ phân kỳ nội họ ở $k > G$, việc chưa triển khai thực nghiệm đối chứng trực tiếp với các phương pháp cố định ngân sách $k$ như chọn ngẫu nhiên, 1-medoid-per-subfamily hoặc CD-HIT/MMseqs2 tinh chỉnh nhị phân ngưỡng tương đồng là một giới hạn phạm vi thực nghiệm cần được bổ sung.

\subsection{Khoảng cách giữa cận xấp xỉ lý thuyết và giải thuật heuristic thực nghiệm, ảnh hưởng của ngưỡng ngân sách tối đa và phương sai xáo trộn}
\label{subsec:lim_static_hyperparams}

Hạn chế thứ hai liên quan đến khoảng cách giữa bảo chứng toán học lý thuyết và giải thuật thực thi trong thực nghiệm. Cận xấp xỉ lý thuyết $(1 - 1/e) \approx 63{,}2\%$ của định lý Nemhauser~\cite{nemhauser1978} áp dụng chặt chẽ cho giải thuật Standard Greedy tối ưu hóa hàm mục tiêu nội tại $F_\alpha(S)$. Để giảm chi phí tính toán trên tập dữ liệu lớn, quy trình thực nghiệm sử dụng cơ chế cắt tỉa cứng $\theta = 0{,}90$ kết hợp lấy top-$k$ theo tần suất đồng thuận qua $1.000$ vòng xáo trộn và đánh giá bằng công cụ Live BLASTp độc lập ($E \le 10^{-7}$). Đây là một giải thuật xấp xỉ heuristic thực nghiệm, không kế thừa nguyên vẹn cận bảo đảm xấu nhất $(1 - 1/e)$. Ngoài ra, ma trận $\mathbf{S}_{\text{BLAST}} = \text{pident}/100$ kế thừa từ bài báo gốc chưa tích hợp tỷ lệ bao phủ chiều dài so khớp (qcovs), khiến hai protein chỉ tương đồng ở một motif xúc tác ngắn vẫn có thể kích hoạt điều kiện cắt tỉa.

Bên cạnh đó, việc trích xuất tập đại diện $S_k$ thông qua tiêu chí sắp xếp tần suất $C(v)$ trên toàn bộ $k_{\max}$ bước khiến danh sách đại diện phụ thuộc vào ngưỡng ngân sách tối đa $k_{\max}$ của từng vòng xáo trộn: tại lát cắt $k = 115$, việc mở rộng $k_{\max}$ từ $115$ lên $380$ làm thay đổi thứ hạng đồng thuận và giảm nhẹ chỉ số Rare Sensitivity từ $59{,}11\%$ xuống $50{,}22\%$. Cuối cùng, toàn bộ kết quả đo đạc trong nghiên cứu hiện được báo cáo trên một danh sách đồng thuận duy nhất khởi tạo từ $\text{seed} = 42$ mà chưa xuất độ lệch chuẩn hay khoảng tin cậy ($\text{mean} \pm \text{std}$) qua nhiều lượt chạy xáo trộn độc lập, đòi hỏi sự cẩn trọng khi diễn giải độ ổn định của các con số thống kê.

\subsection{Hiện tượng bão hòa ma trận tương đồng đặc trên ESM-2, đặc tính bậc thang của thước đo và phạm vi kiểm chứng trên một siêu họ}
\label{subsec:lim_dense_matrix_saturation}

Hạn chế thứ ba xuất phát từ bản chất toán học của ma trận tương đồng đặc trong không gian biểu diễn vector nhúng ESM-2. Kỹ thuật dịch tâm Mean-Centering với $\beta = 0{,}70$ giúp kéo giãn phân bố khoảng cách, song ma trận $\mathbf{S}_{\text{ESM}}^{(\beta=0{,}70)}$ vẫn duy trì tính chất đặc với $99{,}64\%$ phần tử dương.

Khi giải thuật tham lam vận hành ở vùng ngân sách mở rộng ($k \ge 115$), tập nghiệm $S_k$ tích lũy nhiều chuỗi đại diện khiến giá trị bao phủ nền $m_v^{(t)} = \max_{s \in S_t} \mathbf{S}_{\text{ESM}}(s, v)$ trên các chuỗi thuộc phân họ hiếm tăng lên đáng kể do độ tương đồng chéo ngoài đường chéo. Hệ quả là độ tăng cận biên còn lại của các ứng viên nhóm hiếm bị thu hẹp mạnh so với tổng tích lũy độ tăng cận biên của một ứng viên thuộc Subfamily lớn trên hàng trăm phần tử trong nhóm, làm chậm tốc độ bao phủ các phân họ hiếm còn lại so với ma trận thưa BLASTp. Nhận định này hiện là một phân tích giải thích cơ chế dựa trên số liệu phân bổ vị trí đại diện của thử nghiệm 5; các thực nghiệm kiểm chứng cô lập như cắt ngưỡng thưa hóa $\tau$ hay quét tham số $\beta \in [0, 1]$ vẫn cần được tiến hành một cách hệ thống trong các nghiên cứu tiếp theo.

Mặt khác, các chỉ số đo lường cũng mang những đặc tính biên cần lưu ý: chỉ số Worst-Group Sensitivity có đặc tính bậc thang và duy trì bằng $0{,}0\%$ chừng nào còn ít nhất 1 phân họ chưa có đại diện; trong $45$ Subfamily hiếm có tới $18$ phân họ đơn độc ($n_g = 1$) mang tính chất nhị phân $\{0\%, 100\%\}$; và chỉ số Live BLASTp Sensitivity của Subfamily lớn nhất A01.009 bão hòa ở trần $100{,}0\%$ ngay từ $k = 115$. Cuối cùng, toàn bộ các kiểm chứng thực nghiệm của đồ án mới tập trung chuyên sâu trên một siêu họ enzyme duy nhất là MEROPS Family A1; việc kiểm chứng tính tổng quát hóa trên các siêu họ khác như Thioesterases hay Carboxylesterase là yêu cầu quan trọng để khẳng định tính phổ quát của phương pháp luận.

\section{Định hướng nghiên cứu và phát triển tiếp theo}
\label{sec:conc_future_work}

Từ ba ranh giới hạn chế đã được phân tích, đồ án xác lập ba định hướng nghiên cứu và phát triển tiếp theo nhằm nâng tầm khung phương pháp luận RareProtSelect theo hướng thích nghi và tổng quát hơn.

\subsection{Ước lượng mật độ cục bộ và học bán giám sát cho khối chuỗi chưa phân loại}
\label{subsec:fut_upa_semisupervised}

Để giải quyết hạn chế về sự phụ thuộc vào nhãn phân loại sẵn có, hướng phát triển trực tiếp là chuyển đổi cơ chế tính trọng số từ cấp độ phân nhóm rời rạc sang cấp độ mật độ lân cận liên tục. Cụ thể, quy mô mẫu phân họ $n_{g(v)}$ có thể được thay thế bằng hàm ước lượng mật độ cục bộ trên đồ thị lân cận $k$-NN hoặc ước lượng mật độ nhân:
\begin{equation}
    \hat{\rho}(v) = \sum_{u \in \mathcal{N}_K(v)} \mathbf{S}(u, v), \quad \tilde{w}_v^{(\alpha)} = \left(\hat{\rho}(v)\right)^{-\alpha}
    \label{eq:fut_density_weight}
\end{equation}
trong đó $\mathcal{N}_K(v)$ là tập $K$ láng giềng gần nhất của chuỗi $v$. Công thức trọng số liên tục $\tilde{w}_v^{(\alpha)}$ cho phép thuật toán tự động gán trọng số cao cho các chuỗi protein nằm ở các vùng mật độ thưa thớt mà không cần biết trước nhãn phân loại.

Bên cạnh đó, đối với khối $9.187$ chuỗi chưa phân loại [A01.UPA], hướng tiếp cận khả thi là áp dụng các giải thuật lan truyền nhãn bán giám sát hoặc gom cụm phân cấp trên không gian vector nhúng ESM-2 đã qua dịch tâm. Bằng cách sử dụng $4.120$ chuỗi đã có nhãn chính thức làm điểm tựa lan truyền, hệ thống có thể gán nhãn phân họ dự kiến kèm độ tin cậy xác suất cho từng chuỗi [A01.UPA], đồng thời nhận diện các cụm chuỗi cô lập mới như những ứng viên tiềm năng cho các Subfamily mới phát hiện.

\subsection{Thiết kế trọng số thích nghi theo độ phân kỳ nội nhóm và tối ưu hóa dưới ràng buộc Matroid}
\label{subsec:fut_adaptive_matroid}

Để khắc phục giới hạn của siêu tham số tĩnh, một hướng cải tiến tiềm năng là tích hợp độ phân kỳ chuỗi nội nhóm vào công thức trọng số thông qua chỉ số phân tán nội phân họ $D_g$ đo lường khoảng cách trung bình giữa các cá thể trong cùng một Subfamily:
\begin{equation}
    D_g = 1 - \frac{1}{|V_g|^2} \sum_{u \in V_g} \sum_{v \in V_g} \mathbf{S}(u, v)
    \label{eq:fut_dispersion}
\end{equation}
Trọng số của chuỗi khi đó có thể được thiết kế phụ thuộc đồng thời vào quy mô $n_g$ và độ phân tán $D_g$, cho phép tăng cường đại diện cho các phân họ vừa hiếm vừa có độ đa dạng nội nhóm cao.

Đồng thời, bài toán tối ưu hóa có thể được mở rộng từ dạng ràng buộc lực lượng đơn giản sang bài toán tối ưu hóa submodular dưới ràng buộc cấu trúc Partition Matroid:
\begin{equation}
    \max_{S \subseteq V} F_\alpha(S) \quad \text{thỏa mãn} \quad |S \cap V_g| \le \overline{k}_g, \quad \forall g \in \{1, \dots, G\}
    \label{eq:fut_matroid_constraint}
\end{equation}
trong đó $\overline{k}_g$ khống chế trần số lượng đại diện cho từng phân họ nhằm kiểm soát trực tiếp sự dư thừa ở các nhóm lớn. Đối với yêu cầu bảo đảm số lượng đại diện tối thiểu $\underline{k}_g \ge 1$ cho các phân họ hiếm, bài toán có thể được giải quyết bằng kỹ thuật gieo mầm khởi tạo (warm-start seeding) hoặc tích hợp số hạng phạt vi phạm. Về mặt giải thuật, cơ chế cắt tỉa cứng hiện tại có thể được thay thế bằng giải thuật tham lam tăng tốc sử dụng hàng đợi ưu tiên cực đại Lazy Greedy dựa trên cơ chế đánh giá lười CELF~\cite{minoux1978,leskovec2007,krause2014}, vừa tránh rủi ro bỏ sót ứng viên tiềm năng vừa duy trì bảo đảm cận xấp xỉ lý thuyết $1/2$ giá trị tối ưu toàn cục cho bài toán tối ưu hóa submodular đơn điệu dưới ràng buộc matroid.

\subsection{Xây dựng nhân tương đồng lai ghép đa phương thức và mở rộng kiểm chuẩn liên siêu họ}
\label{subsec:fut_hybrid_kernel}

Nhằm kết hợp ưu thế của ma trận tương đồng đặc ở ngân sách thấp và ma trận thưa ở ngân sách cao, hướng phát triển thứ ba là xây dựng ma trận tương đồng lai ghép đa phương thức kết hợp cơ chế cắt ngưỡng để làm thưa hóa ma trận ESM-2:
\begin{equation}
    \mathbf{S}_{\text{hybrid}}(u, v) = \gamma \, \mathbf{S}_{\text{BLAST}}(u, v) + (1 - \gamma) \, \mathbb{I}\!\left[\mathbf{S}_{\text{ESM}}^{(\beta)}(u, v) \ge \tau\right] \cdot \mathbf{S}_{\text{ESM}}^{(\beta)}(u, v)
    \label{eq:fut_hybrid_kernel}
\end{equation}
trong đó $\gamma \in [0, 1]$ là hệ số cân bằng giữa hai không gian biểu diễn và $\tau > 0$ là ngưỡng cắt độ tương đồng tối thiểu nhằm gán về $0$ các giá trị tương đồng yếu mang tính nhiễu nền. Ma trận lai ghép này có thể tích hợp thêm chỉ số tương đồng cấu trúc không gian ba chiều TM-score từ các mô hình dự đoán cấu trúc như AlphaFold2 hoặc ESMFold, giúp kết hợp hài hòa giữa độ chính xác phân họ của so khớp chuỗi, ngữ nghĩa tiến hóa của mô hình ngôn ngữ và hình học không gian phân tử.

Cuối cùng, các nghiên cứu trong tương lai có thể mở rộng kiểm chứng thực nghiệm khung phương pháp luận RareProtSelect trên những siêu họ enzyme khác có cấu trúc phân loại đa dạng, chẳng hạn như siêu họ Thioesterases trên cơ sở dữ liệu ThYme hoặc siêu họ Carboxylesterase trên cơ sở dữ liệu ESTHER. Bên cạnh đó, việc phát triển và đóng gói toàn bộ thuật toán cùng các cấu hình trọng số thành một thư viện phần mềm mã nguồn mở với giao diện lập trình trực quan cũng là một hướng phát triển thiết thực, giúp cộng đồng nghiên cứu tin sinh học và sinh học tiến hóa thuận tiện ứng dụng vào các bài toán phân tích dữ liệu protein quy mô lớn.